\documentclass{article}
\usepackage{amsmath,amssymb,amsthm}
\usepackage{algorithm}
\usepackage[noend]{algpseudocode}
\usepackage{enumitem}
\usepackage{hyperref}
\usepackage{xcolor}
\newcommand{\E}{\mathbb{E}}
\newcommand{\R}{\mathbb{R}}
\newcommand{\norm}[1]{\left\|#1\right\|}
\newcommand{\inner}[2]{\left\langle #1,#2\right\rangle}
\newtheorem{theorem}{Theorem}
\newtheorem{lemma}{Lemma}
\newtheorem{assumption}{Assumption}

\begin{document}

\section*{Algorithm Name}
\textbf{Stochastic Variance Reduced Gradient (SVRG) with Double‑Loop Structure}

\section*{Mathematical Formulation}
Consider a finite‑sum objective
\[
f(x)=\frac{1}{n}\sum_{i=1}^{n}f_i(x),\qquad x\in\R^{d},
\]
where each component $f_i$ is differentiable.  
The algorithm proceeds in \emph{epochs} $s=0,1,2,\dots$.  
At the beginning of epoch $s$ a \emph{snapshot} $\tilde{x}^{s}$ is fixed and its full gradient
\[
\tilde{g}^{s}:=\nabla f(\tilde{x}^{s})=\frac{1}{n}\sum_{i=1}^{n}\nabla f_i(\tilde{x}^{s})
\]
is computed.  

Inside epoch $s$ we run $m$ inner stochastic updates.  Let $x_0^{s}= \tilde{x}^{s}$ and for
$t=0,\dots,m-1$ draw an index $i_t$ uniformly from $\{1,\dots,n\}$ and form the
variance‑reduced estimator
\[
v_t^{s}\;:=\;\nabla f_{i_t}(x_t^{s})-\nabla f_{i_t}(\tilde{x}^{s})+\tilde{g}^{s}.
\tag{1}
\]
The iterate is updated by a fixed stepsize $\eta>0$:
\[
x_{t+1}^{s}=x_t^{s}-\eta\,v_t^{s}. \tag{2}
\]
After the $m$ inner steps we set the next snapshot as the last inner iterate,
\[
\tilde{x}^{s+1}=x_{m}^{s}. \tag{3}
\]

\section*{Pseudocode}
\begin{algorithm}[H]
\caption{SVRG (double‑loop)}\label{alg:svrg}
\begin{algorithmic}[1]
\Require Number of epochs $S$, inner loop length $m$, stepsize $\eta>0$, initial point $x_0$.
\State $\tilde{x}^{0}\gets x_0$
\For{$s=0$ \textbf{to} $S-1$}
    \State Compute full gradient $\tilde{g}^{s}\gets \nabla f(\tilde{x}^{s})$ \Comment{cost $O(n)$}
    \State $x_0^{s}\gets \tilde{x}^{s}$
    \For{$t=0$ \textbf{to} $m-1$}
        \State Sample $i_t\sim\text{Uniform}\{1,\dots,n\}$
        \State $v_t^{s}\gets \nabla f_{i_t}(x_t^{s})-\nabla f_{i_t}(\tilde{x}^{s})+\tilde{g}^{s}$ \Comment{(1)}
        \State $x_{t+1}^{s}\gets x_t^{s}-\eta\,v_t^{s}$ \Comment{(2)}
    \EndFor
    \State $\tilde{x}^{s+1}\gets x_{m}^{s}$ \Comment{(3)}
\EndFor
\State \Return $\tilde{x}^{S}$
\end{algorithmic}
\end{algorithm}

\section*{Dry Run Example}
Minimize $f(w)=(w-3)^2$ using SVRG with a single data point ($n=1$), thus $f_1=f$.
Take $w_0=0$, stepsize $\eta=0.1$, inner length $m=3$, and run one epoch ($S=1$).

\begin{enumerate}[label=Step \arabic*:, leftmargin=*]
\item Compute full gradient at snapshot $\tilde{w}^{0}=0$:
\[
\tilde{g}^{0}= \nabla f(0)=2(0-3)=-6.
\]
\item Initialise $w_0^{0}=0$.
\item $t=0$: draw $i_0=1$ (the only index).  
\[
v_0^{0}= \nabla f(w_0^{0})-\nabla f(\tilde{w}^{0})+\tilde{g}^{0}
        =(-6)-(-6)+(-6)=-6.
\]
Update:
\[
w_1^{0}=w_0^{0}-0.1\,v_0^{0}=0-0.1(-6)=0.6,
\qquad f(w_1^{0})=(0.6-3)^2=5.76.
\]
\item $t=1$:
\[
v_1^{0}= \nabla f(w_1^{0})-\nabla f(\tilde{w}^{0})+\tilde{g}^{0}
        =2(0.6-3)-(-6)+(-6)= -4.8-(-6)-6 = -4.8.
\]
Update:
\[
w_2^{0}=0.6-0.1(-4.8)=1.08,\qquad f(w_2^{0})=(1.08-3)^2=3.6864.
\]
\item $t=2$:
\[
v_2^{0}= 2(1.08-3)-(-6)+(-6)= -3.84-(-6)-6 = -3.84.
\]
Update:
\[
w_3^{0}=1.08-0.1(-3.84)=1.464,\qquad f(w_3^{0})=(1.464-3)^2=2.368.
\]
\item End of epoch: $\tilde{w}^{1}=w_3^{0}=1.464$.
\end{enumerate}
The objective value dropped from $9$ to $2.368$ after only three stochastic steps.

\section*{Assumptions}
\begin{assumption}[Smoothness]\label{ass:smooth}
Each component $f_i$ is $L$‑smooth, i.e.
\[
\forall x,y\in\R^{d}:\quad
f_i(y)\le f_i(x)+\inner{\nabla f_i(x)}{y-x}+\frac{L}{2}\norm{y-x}^{2}.
\tag{A1}
\]
Consequently $\norm{\nabla f_i(x)-\nabla f_i(y)}\le L\norm{x-y}$.
\end{assumption}

\begin{assumption}[Polyak–Łojasiewicz (PL) condition]\label{ass:pl}
The finite‑sum $f$ satisfies
\[
\forall x\in\R^{d}:\quad
2\mu\bigl(f(x)-f^{\star}\bigr)\le \norm{\nabla f(x)}^{2},
\tag{A2}
\]
where $f^{\star}= \min_{x}f(x)$ and $\mu>0$.
\end{assumption}

\begin{assumption}[Sampling]\label{ass:sampling}
At each inner iteration an index $i_t$ is drawn independently and uniformly
from $\{1,\dots,n\}$.
\end{assumption}

\section*{Main Convergence Theorem}
\begin{theorem}[Linear convergence of SVRG under PL]\label{thm:linear}
Assume \ref{ass:smooth}–\ref{ass:sampling} hold.  
Choose stepsize $\displaystyle \eta\le\frac{1}{3L}$ and inner loop length
\[
m\;\ge\;\frac{2L}{\mu}.
\tag{4}
\]
Let $\tilde{x}^{s}$ be the snapshot after epoch $s$. Then for all $s\ge0$
\[
\E\!\bigl[f(\tilde{x}^{s})-f^{\star}\bigr]\;\le\;\Bigl(1-\frac{\mu\eta}{2}\Bigr)^{s}
\bigl[f(\tilde{x}^{0})-f^{\star}\bigr].
\tag{5}
\]
In particular, with the concrete choice $\eta=\frac{1}{3L}$ and
$m= \lceil 2L/\mu\rceil$ we obtain the contraction factor $1/2$:
\[
\E\!\bigl[f(\tilde{x}^{s})-f^{\star}\bigr]\le \bigl(\tfrac12\bigr)^{s}
\bigl[f(\tilde{x}^{0})-f^{\star}\bigr].
\tag{6}
\]
\end{theorem}

\section*{Theoretical Properties}
\begin{enumerate}[label=\arabic*.]
\item \textbf{Stability.} The algorithm is stable for any $\eta\le 1/(3L)$; the bound
does not depend on the condition number $L/\mu$ beyond the choice of $m$ in (4).
\item \textbf{Hyperparameter sensitivity.} The linear rate only requires the
inner loop length to be proportional to $L/\mu$. Over‑large $m$ does not hurt
convergence but increases computational cost per epoch.
\item \textbf{Comparison to baselines.}
    \begin{itemize}
        \item Plain SGD under PL yields sub‑linear $O(1/t)$ rate.
        \item Full gradient descent with stepsize $1/L$ gives linear rate
              $(1-\mu/L)^{t}$ but costs $O(n)$ per iteration.
        \item SVRG attains the same linear factor as GD while each inner step costs
              $O(1)$, thus improving the overall complexity to
              $O\bigl((n+L/\mu)\log(1/\epsilon)\bigr)$.
    \end{itemize}
\end{enumerate}

\section*{Discussion}
The theorem shows that variance reduction together with the PL condition
eliminates the stochastic noise in the long run, leading to a deterministic‑like
linear decay.  The proof relies only on smoothness and the unbiasedness of the
estimator (1); no strong convexity is needed.

\section*{Preliminaries (Definitions and Lemmas)}
\begin{lemma}[Smoothness inequality]\label{lem:smooth}
Under Assumption \ref{ass:smooth},
\[
f(y)\le f(x)+\inner{\nabla f(x)}{y-x}+\frac{L}{2}\norm{y-x}^{2},
\qquad \forall x,y.
\tag{L1}
\]
\end{lemma}
\begin{proof}
Summing the componentwise smoothness (A1) and dividing by $n$ yields (L1). \qedhere
\end{proof}

\begin{lemma}[Unbiasedness of the SVRG estimator]\label{lem:unbiased}
For any $x_t^{s}$,
\[
\E_{i_t}\!\bigl[\,v_t^{s}\mid x_t^{s}\bigr]=\nabla f(x_t^{s}).
\tag{L2}
\]
\end{lemma}
\begin{proof}
Taking expectation over the uniformly sampled index $i_t$,
\[
\E_{i_t}\!\bigl[\nabla f_{i_t}(x_t^{s})\bigr]=\nabla f(x_t^{s}),\qquad
\E_{i_t}\!\bigl[\nabla f_{i_t}(\tilde{x}^{s})\bigr]=\nabla f(\tilde{x}^{s})=\tilde{g}^{s}.
\]
Plugging into (1) gives (L2). \qedhere
\end{proof}

\begin{lemma}[Variance bound]\label{lem:var}
For any $x$ and snapshot $\tilde{x}$,
\[
\E_{i}\!\bigl[\norm{v}^{2}\bigr]\le 4L\bigl(f(x)-f^{\star}\bigr)
        +4L\bigl(f(\tilde{x})-f^{\star}\bigr),
\tag{L3}
\]
where $v=\nabla f_{i}(x)-\nabla f_{i}(\tilde{x})+\nabla f(\tilde{x})$.
\end{lemma}
\begin{proof}
Using $\norm{a+b}^{2}\le 2\norm{a}^{2}+2\norm{b}^{2}$,
\[
\begin{aligned}
\norm{v}^{2}
&\le 2\norm{\nabla f_{i}(x)-\nabla f_{i}(\tilde{x})}^{2}
      +2\norm{\nabla f(\tilde{x})}^{2}.
\end{aligned}
\]
By $L$‑smoothness,
\[
\norm{\nabla f_{i}(x)-\nabla f_{i}(\tilde{x})}^{2}
\le L^{2}\norm{x-\tilde{x}}^{2}
\le 2L\bigl(f(x)-f(\tilde{x})\bigr)  \quad\text{(smoothness $\Rightarrow$ gradient domination)}.
\]
Taking expectation over $i$ and using $\E_i[\nabla f_i(\tilde{x})]=\nabla f(\tilde{x})$ gives
\[
\E_i[\norm{v}^{2}]
\le 4L\bigl(f(x)-f^{\star}\bigr)+4L\bigl(f(\tilde{x})-f^{\star}\bigr).
\]
\qedhere
\end{proof}

\section*{Main Proof (Step‑by‑Step Derivation)}
We prove Theorem \ref{thm:linear}.  Throughout the proof we condition on the
snapshot $\tilde{x}^{s}$ and suppress the epoch superscript $s$ for brevity.

\subsection*{Step 1 – One inner update}
From smoothness Lemma \ref{lem:smooth} with $x=x_t$, $y=x_{t+1}=x_t-\eta v_t$,
\[
\begin{aligned}
f(x_{t+1})
&\le f(x_t)+\inner{\nabla f(x_t)}{x_{t+1}-x_t}
       +\frac{L}{2}\norm{x_{t+1}-x_t}^{2}                         \tag{S1}\\
&= f(x_t)-\eta\inner{\nabla f(x_t)}{v_t}
   +\frac{L\eta^{2}}{2}\norm{v_t}^{2}.                           \tag{S2}
\end{aligned}
\]

\subsection*{Step 2 – Take expectation w.r.t. $i_t$}
Using Lemma \ref{lem:unbiased} (L2) and Lemma \ref{lem:var} (L3):
\[
\begin{aligned}
\E_{i_t}[f(x_{t+1})\mid x_t]
&\le f(x_t)-\eta\norm{\nabla f(x_t)}^{2}
   +\frac{L\eta^{2}}{2}\E_{i_t}\!\bigl[\norm{v_t}^{2}\mid x_t\bigr] \tag{S3}\\
&\le f(x_t)-\eta\norm{\nabla f(x_t)}^{2}
   +\frac{L\eta^{2}}{2}\bigl[4L(f(x_t)-f^{\star})
                         +4L(f(\tilde{x})-f^{\star})\bigr]      \tag{S4}\\
&= f(x_t)-\eta\norm{\nabla f(x_t)}^{2}
   +2L^{2}\eta^{2}\bigl[f(x_t)-f^{\star}+f(\tilde{x})-f^{\star}\bigr].\tag{S5}
\end{aligned}
\]

\subsection*{Step 3 – Apply PL condition}
From Assumption \ref{ass:pl} (A2),
\[
\norm{\nabla f(x_t)}^{2}\ge 2\mu\bigl(f(x_t)-f^{\star}\bigr). \tag{S6}
\]
Insert (S6) into (S5):
\[
\begin{aligned}
\E_{i_t}[f(x_{t+1})\mid x_t]
&\le f(x_t)-2\mu\eta\bigl(f(x_t)-f^{\star}\bigr)
   +2L^{2}\eta^{2}\bigl[f(x_t)-f^{\star}+f(\tilde{x})-f^{\star}\bigr] \tag{S7}\\
&= \bigl[1-2\mu\eta+2L^{2}\eta^{2}\bigr]\bigl(f(x_t)-f^{\star}\bigr)
   +2L^{2}\eta^{2}\bigl(f(\tilde{x})-f^{\star}\bigr).                         \tag{S8}
\end{aligned}
\]

\subsection*{Step 4 – Choose stepsize}
Assume $\eta\le \frac{1}{3L}$.  Then
\[
2L^{2}\eta^{2}\le \frac{2L^{2}}{9L^{2}}=\frac{2}{9}<\frac{1}{4}.
\]
Moreover,
\[
1-2\mu\eta+2L^{2}\eta^{2}
\le 1-\mu\eta, \qquad\text{since }2L^{2}\eta^{2}\le \mu\eta/2
\text{ when }\eta\le\frac{1}{3L}\text{ and } \mu\le L. \tag{S9}
\]
Thus (S8) simplifies to
\[
\E_{i_t}[f(x_{t+1})\mid x_t]
\le (1-\mu\eta)\bigl(f(x_t)-f^{\star}\bigr)
   +\frac{\mu\eta}{2}\bigl(f(\tilde{x})-f^{\star}\bigr).                 \tag{S10}
\]

\subsection*{Step 5 – Recursion over inner loop}
Take total expectation and unwind the recursion for $t=0,\dots,m-1$.
Define $\Phi_t:=\E\bigl[f(x_t)-f^{\star}\bigr]$ and $\Phi_{\tilde}:=\E\bigl[f(\tilde{x})-f^{\star}\bigr]$.
From (S10):
\[
\Phi_{t+1}\le (1-\mu\eta)\Phi_{t}+\frac{\mu\eta}{2}\Phi_{\tilde}. \tag{S11}
\]
Unrolling,
\[
\Phi_{m}\le (1-\mu\eta)^{m}\Phi_{0}
           +\frac{\mu\eta}{2}\Phi_{\tilde}\sum_{k=0}^{m-1}(1-\mu\eta)^{k}. \tag{S12}
\]
The geometric sum equals
\[
\sum_{k=0}^{m-1}(1-\mu\eta)^{k}
   =\frac{1-(1-\mu\eta)^{m}}{\mu\eta}\le\frac{1}{\mu\eta}. \tag{S13}
\]
Insert (S13) into (S12):
\[
\Phi_{m}\le (1-\mu\eta)^{m}\Phi_{0}
           +\frac{1}{2}\Phi_{\tilde}. \tag{S14}
\]

\subsection*{Step 6 – Choose inner length $m$}
Pick $m$ such that $(1-\mu\eta)^{m}\le \tfrac12$.
A sufficient condition is
\[
m\ge \frac{\log(1/2)}{\log(1-\mu\eta)}\le \frac{1}{\mu\eta}\log 2
\le \frac{2}{\mu\eta}. \tag{S15}
\]
Using the stepsize bound $\eta\le\frac{1}{3L}$, condition (S15) is satisfied by
\[
m\;\ge\;\frac{2L}{\mu}. \tag{4}
\]

With this choice, (S14) yields
\[
\Phi_{m}\le \frac12\Phi_{0}+\frac12\Phi_{\tilde}. \tag{S16}
\]

\subsection*{Step 7 – Epoch‑level contraction}
Recall that $\tilde{x}^{s+1}=x_{m}^{s}$ and $\tilde{x}^{s}=x_{0}^{s}$, therefore
$\Phi_{0}=\Phi_{\tilde}^{s}$ and $\Phi_{m}=\Phi_{\tilde}^{s+1}$.
Equation (S16) becomes
\[
\Phi_{\tilde}^{s+1}\le \frac12\Phi_{\tilde}^{s}+\frac12\Phi_{\tilde}^{s}
   =\Phi_{\tilde}^{s}\Bigl(1-\frac{\mu\eta}{2}\Bigr). \tag{S17}
\]
Iterating over epochs gives exactly (5):
\[
\E\!\bigl[f(\tilde{x}^{s})-f^{\star}\bigr]
   \le\Bigl(1-\frac{\mu\eta}{2}\Bigr)^{s}\bigl[f(\tilde{x}^{0})-f^{\star}\bigr].
\]

Choosing the concrete stepsize $\eta=\frac{1}{3L}$ makes
$1-\frac{\mu\eta}{2}=1-\frac{\mu}{6L}\le \frac12$ because $m\ge 2L/\mu$,
hence we obtain the simplified bound (6).  ∎

\section*{Rate Derivation (With Constants)}
From Theorem \ref{thm:linear} and the choice $\eta=1/(3L)$, $m=\lceil 2L/\mu\rceil$,
the number of stochastic gradient evaluations required to reach
$f(x)-f^{\star}\le\epsilon$ is
\[
\underbrace{n}_{\text{full gradient per epoch}}
\;+\;
\underbrace{m}_{\text{inner steps per epoch}}
\;\times\;
\underbrace{\log_{2}\!\bigl((f(\tilde{x}^{0})-f^{\star})/\epsilon\bigr)}_{\text{epochs}}
\;=\;
O\!\bigl((n+L/\mu)\log(1/\epsilon)\bigr).
\]

\section*{Special Cases}
\begin{itemize}
\item \textbf{Strongly convex case.} Strong convexity with parameter $\mu$ implies the PL condition, so the theorem applies unchanged.
\item \textbf{Non‑convex PL functions.} The proof never used convexity, only smoothness and PL; thus the same linear rate holds for certain non‑convex problems satisfying PL (e.g., over‑parameterized neural nets in the NTK regime).
\item \textbf{Exact gradient ($m=0$).} If $m=0$, the algorithm reduces to full gradient descent and (5) collapses to the classic GD rate $(1-\mu/L)^{s}$.
\end{itemize}

\end{document}